\documentclass[conference]{IEEEtran}
\IEEEoverridecommandlockouts

\usepackage{cite}
\usepackage{amsmath,amssymb,amsfonts}
\usepackage{algorithmic}
\usepackage{graphicx}
\usepackage{textcomp}
\usepackage{xcolor}
\usepackage{booktabs}
\usepackage{hyperref}
\usepackage{enumitem}

\def\BibTeX{{\rm B\kern-.05em{\sc i\kern-.025em b}\kern-.08em
    T\kern-.1667em\lower.7ex\hbox{E}\kern-.125emX}}

\begin{document}

\title{Mixture-of-Experts Framework for \\ High-Frequency Mid-Price Movement Prediction}

\author{\IEEEauthorblockN{Authors}
\IEEEauthorblockA{\textit{Pattern Recognition and Machine Learning (PRML)} \\
\textit{Institute}\\
City, Country \\
Email}
}

\maketitle

\begin{abstract}
Predicting short-term mid-price movements in Limit Order Books (LOBs) is 
a profoundly difficult financial engineering problem due to extreme 
non-stationarity, microstructure noise, and severe class imbalance. This 
work introduces a Mixture-of-Experts (MoE) pipeline on the FI-2010 
benchmark dataset combining L2-Regularised Logistic Regression, 
Gradient-Boosted Decision Trees (XGBoost), and a Temporal Multi-Layer 
Perceptron (MLP), orchestrated via a learned Gating Network with a 
20-step signal regime filter. We rigorously investigate the 
\textit{Classification-Profitability Paradox}---demonstrating that the 
highest-accuracy model (MLP, Macro-F1 = 56.2\%) achieves near-zero 
financial return at 1bp cost, while lower-accuracy XGBoost generates 
positive returns by maintaining directional conviction above execution 
thresholds. Across 9 chronological folds, the Gated MoE achieves a mean 
return of $+19.98\%$ at 1bp, with superior risk-adjusted performance 
(Sortino = 0.1224) compared to XGBoost solo (Sortino = 0.0969), though 
the raw return difference is not statistically significant ($p = 0.90$). 
A critical finding is that naive ensemble averaging (equal-weight, 
F1-weighted) collapses returns to near zero by diluting conviction below 
the $\tau_{entry}=0.60$ execution threshold---confirming that a learned 
gating mechanism is necessary, if not sufficient, for financial extraction.
\end{abstract}

\begin{IEEEkeywords}
Mixture of Experts, Limit Order Book, High-Frequency Trading, XGBoost, Neural Networks, Focal Loss, Order Imbalance
\end{IEEEkeywords}

\section{Introduction}
Financial markets have evolved into highly digitized ecosystems dominated by algorithmic trading protocols operating at the millisecond scale. The fundamental data structure dictating price discovery in these modern markets is the Limit Order Book (LOB), a real-time sequential ledger matching standing buy (Bid) and sell (Ask) liquidity. 

The core predictive challenge intrinsic to the LOB is extreme market microstructure noise. Unlike macroscopic asset allocation, HFT signals decay in seconds, and mid-price trajectories often revert due to phantom liquidity or transient order cancellations. Consequently, standard single-model machine learning regressors are notoriously susceptible to overfitting against this localized variance.

To mitigate this, financial literature has widely adopted categorical formulation. Predicting the exact scalar return $R_t$ is abandoned in favor of mapping the state to three discrete classifications: upward movement (Up), downward movement (Down), and no significant movement (Stationary).

However, this transition introduces the "Classification-Profitability Paradox." A standard machine learning model trained via discrete Cross-Entropy loss will structurally gravitate toward the majority class (Stationary). While capturing the Stationary class aggressively yields high systemic accuracy metrics, it paralyzes the agent; the model rarely builds enough conviction to physically place a directional trade. Thus, achieving $90\%$ accuracy yields $0\%$ financial extraction. 

This project solves this paradox by decoupling the classification mechanism from the financial execution mechanism using a Mixture-of-Experts architecture routed through a highly constrained, cost-aware mathematical state machine.

\section{Dataset Configuration and Preprocessing}
We utilize the spatial FI-2010 benchmark LOB dataset \cite{b1}, widely considered the academic standard for microstructure forecasting. The corpus encapsulates 10 consecutive trading days consisting of nearly 4 million tick events, with visibility across the outer 10 price levels of the limit order book.

\subsection{Raw Decimal Precision Extraction}
The FI-2010 benchmark provides arrays normalized via Z-Score. However, utilizing standardized distributions destroys the strict positivity of order volume integers. Subtracting absolute integers mathematically corrupts non-linear ratio engineering (such as calculating relative volume imbalances). 

Therefore, our pipeline bypasses the generalized subsets and forces the extraction of the raw \texttt{NoAuction\_DecPre} (Decimal Precision) matrices. By evaluating physical, absolute quantities, we guarantee mathematically pure fractional inputs for our algorithms.

\subsection{Dimensionality and Feature Engineering}
The generalized dataset provides 144 baseline spatial features (including hierarchical price points, cumulative volumes, and basic localized moving averages). Following exploratory data analysis (EDA), three critical, high-conviction features structurally isolated directional momentum:

\begin{enumerate}
    \item \textbf{Level-1 Spread ($S_t$):} Explicitly mapping the threshold required to instantly cross liquidity.
    \begin{equation}
        S_t = AskPrice_1 - BidPrice_1
    \end{equation}
    
    \item \textbf{Level-1 Market Imbalance ($I_t^{(1)}$):} Normalizing the immediate physical buying pressure against the immediate selling pressure into a bounded $[-1, 1]$ scalar.
    \begin{equation}
        I_t^{(1)} = \frac{BidVol_1 - AskVol_1}{BidVol_1 + AskVol_1}
    \end{equation}
    
    \item \textbf{Level-5 Deep Imbalance ($I_t^{(5)}$):} A holistic discrepancy scalar smoothed across the cumulative outer 5 tiers of the order book, intentionally designed to nullify Level-1 high-frequency spoofing noise.
\end{enumerate}

These additions expanded the parameter tensor horizontally to a dimension of 147. To guarantee gradient stability and model convergence across disjointed statistical bands, the entire 147-dimensional matrix is subjected to an \texttt{sklearn.StandardScaler} internally within the dataset loader immediately prior to inference.

\section{Mathematical Formulation}

\subsection{The Label Smoothing Target}
Following standard HFT classification formulation, the physical percentage change of the mid-price $m_t$ evaluated horizontally over a horizon of $k=10$ subsequent ticks defines the strict target vector $y_t \in \{Up, Stationary, Down\}$, bounded by a static financial threshold scalar $l$:
\begin{equation}
    y_t = 
    \begin{cases} 
      Up & \text{if } \frac{m_{t+k} - m_t}{m_t} > l \\
      Down & \text{if } \frac{m_{t+k} - m_t}{m_t} < -l \\
      Stationary & \text{otherwise} 
    \end{cases}
\end{equation}

\subsection{Mixture of Experts Function}
The foundational architecture maps the input state $X_t$ independently across $N$ disparate classifiers (Experts). A concurrent Gating Network calculates a contextual softmax probability density $G_i(X_t)$ assigning dynamic trust weights to each expert. The final stacked probability tensor $P(X_t)$ is the sum of weighted interpolations:
\begin{equation}
    P(X_t) = \sum_{i=1}^{N} G_i(X_t) E_i(X_t)
\end{equation}

\section{Model Architecture and Hyperparameter Logic}
Our architecture encapsulates three fundamentally distinct classifiers to guarantee uncorrelated prediction variances. 

\subsection{Expert 1: Logistic Regression}
A strict baseline linear classifier governed by an L2-penalty norm. It possesses zero inherent capacity for interaction discovery (such as conditional depth tracking), forcing its weights to heavily exploit our artificially engineered $I_t^{(1)}$ matrices.

\subsection{Expert 2: XGBoost Decision Trees}
A Gradient Boosted architecture exploiting spatial mapping. In tabular data execution, decision trees inherently outperform sequential deep learning by mapping boolean splits aggressively. To prevent chaotic microstructure overfitting, we aggressively restricted the absolute max tree-depth iteratively. XGBoost dynamically parses boundary interactions within localized feature depths without normalizing requirements.

\subsection{Expert 3: Temporal Multi-Layer Perceptron (MLP)}
Unlike decision trees, deep learning vectors digest chronological sequences exceptionally. We implemented a temporal mapping mechanism specifically for the MLP. The 147-dimensional array is shifted over a chronological memory window of $k_{lookback}=5$. The current timestep $(t)$ is horizontally concatenated adjacent to the historical sequence $(t-4, t-3, t-2, t-1)$, projecting the input shape into an ultra-wide dense tensor:
\begin{equation}
    D_{MLP} = 147 \text{ features} \times 5 \text{ frames} = 735
\end{equation}
To counteract the profound class imbalance toward "Stationary" targets, the PyTorch MLP was optimized via Focal Loss ($\gamma=3.0$), which structurally de-weights well-classified (safe) gradients and magnifies error penalties on hard, minority directional vectors:
\begin{equation}
    FL(p_t) = -\sum (1 - p_t)^\gamma \log(p_t)
\end{equation}

\section{Algorithmic Trading State Machine}
The absolute raw probabilities generated by the MoE $P(X_t)$ are inherently meaningless without execution limits. Trading entails transaction slippage and physical holding durations. We constructed a deterministic Signal Layer constrained by localized hysteresis. 

\subsection{Conviction Hysteresis}
Rather than executing on generic maximum likelihoods ($argmax$), trading requires the directional signal to eclipse the stationary null hypothesis by a rigid magnitude offset, $\tau_{entry} = 0.60$.
To prevent continuous "churning" across tight regimes, existing physical positions are heavily shielded by an asymmetric exit protocol ($\tau_{exit} = 0.15$), requiring the mathematical signal to decay entirely.

\subsection{The Regime Filter}
The most financially corrosive market structure is the "Sideways Chop"—stochastic flickering that randomly triggers entry thresholds but fails to execute momentum. 

To eliminate this volatility trap, we implemented a 20-step sliding execution deque tracking the maximal directional indicator $s_t$:
\begin{equation}
    s_t = \max(P_{up}(X_t), P_{down}(X_t))
\end{equation}
All directional execution functions are physically halted unless the sequence demonstrates sustained sequential volume pressure:
\begin{equation}
    \frac{1}{20} \sum_{i=0}^{19} s_{t-i} > 0.55
\end{equation}
This isolated regime architecture definitively prevented over 100+ "noise-based" losing trades per fold subset. 

\section{Experimental Results and Discussion}
The pipeline was evaluated using Anchored Forward Time-Series 
Cross-Validation across 9 chronologically sequential folds, preventing 
data leakage. All metrics are reported as mean $\pm$ std across folds.

\subsection{Classification Performance}
\begin{table}[htbp]
\caption{Classification Metrics (9-Fold Mean $\pm$ Std)}
\begin{center}
\begin{tabular}{lcc}
\toprule
\textbf{Model} & \textbf{Accuracy} & \textbf{Macro-F1} \\
\midrule
LR (Expert 1)          & 0.447 $\pm$ 0.019 & 0.432 $\pm$ 0.014 \\
XGBoost (Expert 2)     & 0.537 $\pm$ 0.028 & 0.502 $\pm$ 0.044 \\
MLP (Expert 3)         & \textbf{0.572 $\pm$ 0.036} & \textbf{0.562 $\pm$ 0.033} \\
Gated MoE              & 0.571 $\pm$ 0.038 & 0.534 $\pm$ 0.064 \\
\bottomrule
\end{tabular}
\end{center}
\end{table}

\subsection{The Classification-Profitability Paradox}
Table~\ref{tab_paradox} directly formalises the central paradox. The MLP 
achieves the highest Macro-F1 but near-zero return, while XGBoost 
sacrifices classification accuracy for directional conviction. The Gated 
MoE achieves the highest risk-adjusted performance.

\begin{table}[htbp]
\caption{Classification vs Financial Performance (1bp, 9-Fold Mean)}
\label{tab_paradox}
\begin{center}
\begin{tabular}{lccc}
\toprule
\textbf{Model} & \textbf{Macro-F1} & \textbf{Return} & \textbf{Sortino} \\
\midrule
MLP (Expert 3)     & \textbf{0.562} & $\approx 0\%$  & low \\
XGBoost (Expert 2) & 0.502          & +20.94\%       & 0.097 \\
\textbf{Gated MoE} & 0.534          & +19.98\%       & \textbf{0.122} \\
\bottomrule
\end{tabular}
\end{center}
\end{table}

The MLP's conservative probability outputs rarely breach 
$\tau_{entry}=0.60$, resulting in very few trades and negative net return 
from transaction costs. XGBoost's tabular split mappings produce 
high-conviction directional outputs that consistently exceed the execution 
threshold. The MoE Gating Network learns to preserve this conviction while 
reducing return volatility.

\subsection{Gating Ablation}
\begin{table}[htbp]
\caption{Gating Ablation: Return at 1bp (9-Fold Mean $\pm$ Std)}
\begin{center}
\begin{tabular}{lc}
\toprule
\textbf{Ensemble Method} & \textbf{Mean Return} \\
\midrule
XGB Solo                & +20.94\% $\pm$ 74.97\% \\
Equal-Weight (1/3 each) & +0.47\% $\pm$ 0.86\%   \\
F1-Weighted Blend       & +0.48\% $\pm$ 0.77\%   \\
\textbf{Gated MoE}      & +19.98\% $\pm$ 66.73\% \\
\bottomrule
\end{tabular}
\end{center}
\end{table}

A critical finding: naive ensemble averaging (equal-weight and F1-weighted) 
collapses financial return to near zero ($\approx$+0.5\%). Averaging the 
probability vectors dilutes directional conviction below the 
$\tau_{entry}=0.60$ execution threshold, preventing virtually all trades 
from executing. The learned Gating Network avoids this collapse by 
learning to preserve XGBoost's conviction structure, achieving returns 
comparable to XGB solo with reduced variance.

\subsection{Transaction Cost Sensitivity}
\begin{table}[htbp]
\caption{Cost Sensitivity Analysis}
\begin{center}
\begin{tabular}{lcc}
\toprule
\textbf{Cost} & \textbf{MoE Return} & \textbf{XGB Return} \\
\midrule
1 bp & +19.98\% $\pm$ 66.73\% & +20.94\% $\pm$ 74.97\% \\
3 bp & +6.11\% $\pm$ 55.13\%  & +6.91\% $\pm$ 62.47\%  \\
\bottomrule
\end{tabular}
\end{center}
\end{table}

Returns decay substantially under modest cost increases, reflecting the 
high-frequency nature of the strategy. At real-world execution costs of 
5--10bp, profitability cannot be guaranteed. These figures should be 
interpreted as optimistic simulation estimates.

\subsection{Statistical Rigor}
A paired t-test across 9 folds yields $t = -0.129$, $p = 0.90$, 
indicating no statistically significant raw return difference between 
MoE and XGBoost. The MoE outperforms XGB in 4 of 9 folds. Median returns 
are negative for both strategies (MoE: $-1.93\%$, XGB: $-1.20\%$), 
indicating that high positive returns in a minority of folds drive the 
positive mean---a characteristic of high-variance trading strategies on 
short horizons. The Spearman correlation between Macro-F1 and financial 
return across models yields $\rho = -0.046$ ($p = 0.82$), providing weak 
but directionally consistent evidence for the paradox hypothesis. Due to 
temporal dependence between folds, all statistical tests are approximate.

\section{Conclusion}
This work demonstrates that LOB mid-price prediction cannot be optimised 
solely via classification accuracy. The core finding is threefold. First, 
the Classification-Profitability Paradox is empirically supported: the 
highest-F1 model (MLP) fails financially, while the lower-F1 XGBoost 
generates positive returns by maintaining execution conviction. Second, 
naive ensemble averaging catastrophically suppresses financial performance 
by diluting probability mass below execution thresholds---making a learned 
Gating Network architecturally necessary. Third, the Gated MoE provides 
evidence of risk-adjusted improvement over single experts (higher Sortino, 
lower return variance) without statistically significant raw return gains 
($p = 0.90$). These findings reinforce that ensemble design for financial 
execution must account for threshold-sensitive execution dynamics, not 
only classification objectives.

\begin{thebibliography}{00}
\bibitem{b1} A. Ntakaris et al., ``Benchmark dataset for mid-price forecasting of limit order book data with machine learning methodology.'' Journal of Forecasting, vol. 37, no. 8, pp. 852-866, 2018.
\bibitem{b2} N. Passalis et al., ``Time-Series Classification Using Neural Bag-of-Features.'' 2017.
\end{thebibliography}

\end{document}
